\documentclass[12pt,a4paper]{article}

\usepackage[utf8]{inputenc}
\usepackage[T1]{fontenc}
\usepackage{amsmath,amssymb}
\usepackage{enumitem}
\usepackage{geometry}
\usepackage{multicol}
\usepackage{xcolor}

\geometry{margin=2.3cm}

\title{Banco de Questões Objetivas -- Aula 14\\
Técnicas de Recuperação de Banco de Dados}
\author{Banco de Dados}
\date{}

\begin{document}

\maketitle

\section*{Conteúdos abordados}

\begin{itemize}
    \item objetivos da recuperação de banco de dados;
    \item restauração típica após falhas;
    \item log do sistema;
    \item ponto de efetivação;
    \item imagens anterior e posterior;
    \item UNDO e REDO;
    \item gravação adiantada de log;
    \item atualização in loco e atualização sombra;
    \item políticas \textit{steal}, \textit{no-steal}, \textit{force} e \textit{no-force};
    \item atualização adiada;
    \item atualização imediata;
    \item NO-UNDO/REDO;
    \item UNDO/REDO;
    \item UNDO/NO-REDO;
    \item NO-UNDO/NO-REDO;
    \item checkpoints;
    \item paginação sombra;
    \item recuperação em ambientes multiusuários;
    \item backup e falhas catastróficas.
\end{itemize}

\section{Questões de múltipla escolha}

\begin{enumerate}[label=\textbf{Q\arabic*.}]

\item O objetivo principal das técnicas de recuperação de banco de dados é:

\begin{enumerate}[label=\alph*)]
    \item aumentar a quantidade de tabelas do sistema.
    \item restaurar o banco para um estado consistente após falhas.
    \item substituir completamente o controle de concorrência.
    \item impedir qualquer execução de transações.
\end{enumerate}

\item As técnicas de recuperação contribuem principalmente para garantir a propriedade ACID de:

\begin{enumerate}[label=\alph*)]
    \item atomicidade.
    \item isolamento apenas.
    \item normalização.
    \item independência lógica.
\end{enumerate}

\item Além da atomicidade, a recuperação também está diretamente relacionada à:

\begin{enumerate}[label=\alph*)]
    \item durabilidade.
    \item projeção.
    \item junção.
    \item dependência funcional.
\end{enumerate}

\item Uma falha de transação pode ocorrer devido a:

\begin{enumerate}[label=\alph*)]
    \item erro lógico, violação de restrição ou aborto explícito.
    \item uso de uma chave primária.
    \item execução de uma consulta correta.
    \item existência de índices.
\end{enumerate}

\item Uma falha de sistema pode ser causada por:

\begin{enumerate}[label=\alph*)]
    \item queda de energia ou falha do sistema operacional.
    \item uma consulta com \texttt{ORDER BY}.
    \item uma relação em BCNF.
    \item um atributo \texttt{NOT NULL}.
\end{enumerate}

\item Em uma restauração típica, o sistema deve:

\begin{enumerate}[label=\alph*)]
    \item analisar o estado das transações e restaurar os efeitos corretos.
    \item remover todas as tabelas.
    \item apagar obrigatoriamente o log.
    \item confirmar todas as transações incompletas.
\end{enumerate}

\item O log do sistema é utilizado para:

\begin{enumerate}[label=\alph*)]
    \item registrar informações sobre a execução e atualização das transações.
    \item armazenar somente consultas \texttt{SELECT}.
    \item substituir o banco de dados.
    \item controlar apenas permissões de usuários.
\end{enumerate}

\item Um registro de log associado a uma atualização normalmente contém:

\begin{enumerate}[label=\alph*)]
    \item identificação da transação, item alterado e valores anterior e posterior.
    \item apenas o nome do banco.
    \item somente o horário de início do SGBD.
    \item somente o comando \texttt{COMMIT}.
\end{enumerate}

\item A imagem anterior de um item, também chamada de BFIM, representa:

\begin{enumerate}[label=\alph*)]
    \item o valor do item antes da atualização.
    \item o valor do item após a atualização.
    \item a chave primária da tabela.
    \item o nome da transação.
\end{enumerate}

\item A imagem posterior de um item, também chamada de AFIM, representa:

\begin{enumerate}[label=\alph*)]
    \item o valor do item antes da atualização.
    \item o valor do item depois da atualização.
    \item o estado inicial do servidor.
    \item o identificador do log.
\end{enumerate}

\item A operação UNDO utiliza normalmente:

\begin{enumerate}[label=\alph*)]
    \item a imagem anterior do item.
    \item a imagem posterior do item.
    \item apenas o identificador da tabela.
    \item apenas o comando \texttt{COMMIT}.
\end{enumerate}

\item A operação REDO utiliza normalmente:

\begin{enumerate}[label=\alph*)]
    \item a imagem anterior.
    \item a imagem posterior.
    \item uma chave estrangeira.
    \item o valor nulo.
\end{enumerate}

\item O propósito do UNDO é:

\begin{enumerate}[label=\alph*)]
    \item desfazer os efeitos de uma transação incompleta ou abortada.
    \item repetir uma transação confirmada.
    \item criar um backup.
    \item apagar todo o histórico.
\end{enumerate}

\item O propósito do REDO é:

\begin{enumerate}[label=\alph*)]
    \item refazer atualizações de uma transação confirmada que podem não ter chegado ao disco.
    \item desfazer uma transação abortada.
    \item cancelar todos os commits.
    \item impedir uso do log.
\end{enumerate}

\item O ponto de efetivação de uma transação corresponde ao momento em que:

\begin{enumerate}[label=\alph*)]
    \item a transação é considerada confirmada.
    \item a transação inicia sua primeira leitura.
    \item o banco é apagado.
    \item ocorre necessariamente uma falha.
\end{enumerate}

\item Uma transação incompleta no momento da falha geralmente:

\begin{enumerate}[label=\alph*)]
    \item não deve ter seus efeitos considerados definitivamente confirmados.
    \item deve sempre receber \texttt{COMMIT}.
    \item deve ser tratada como concluída.
    \item não precisa ser analisada.
\end{enumerate}

\item A gravação adiantada de log é conhecida pela sigla:

\begin{enumerate}[label=\alph*)]
    \item WAL.
    \item DDL.
    \item JDBC.
    \item BCNF.
\end{enumerate}

\item A regra básica de \textit{Write-Ahead Logging} determina que:

\begin{enumerate}[label=\alph*)]
    \item o registro de log correspondente deve chegar ao disco antes da página alterada.
    \item a página deve ser escrita antes do log.
    \item o log só deve ser criado após uma falha.
    \item as transações não devem usar \texttt{COMMIT}.
\end{enumerate}

\item Antes de confirmar uma transação, o WAL também exige normalmente que:

\begin{enumerate}[label=\alph*)]
    \item seus registros de log relevantes estejam persistidos.
    \item todas as tabelas sejam bloqueadas permanentemente.
    \item o banco seja reiniciado.
    \item todos os buffers sejam descartados.
\end{enumerate}

\item Uma atualização in loco ocorre quando:

\begin{enumerate}[label=\alph*)]
    \item o novo valor substitui o valor antigo na mesma localização do banco.
    \item uma cópia totalmente separada do banco é criada.
    \item nenhuma página é modificada.
    \item apenas o log é atualizado.
\end{enumerate}

\item Na técnica de atualização sombra:

\begin{enumerate}[label=\alph*)]
    \item as modificações são feitas em páginas alternativas, preservando a versão anterior.
    \item toda atualização sobrescreve imediatamente a página original.
    \item o log é obrigatoriamente o único mecanismo existente.
    \item não existem páginas em disco.
\end{enumerate}

\item Uma vantagem da atualização sombra é:

\begin{enumerate}[label=\alph*)]
    \item possibilitar recuperação sem UNDO e sem REDO em sua forma básica.
    \item exigir obrigatoriamente UNDO e REDO.
    \item impedir qualquer commit.
    \item eliminar o armazenamento em disco.
\end{enumerate}

\item Uma desvantagem da atualização sombra é:

\begin{enumerate}[label=\alph*)]
    \item custo de gerenciamento e possível fragmentação.
    \item impossibilidade de manter uma versão antiga.
    \item necessidade obrigatória de UNDO.
    \item impossibilidade de usar ponteiros.
\end{enumerate}

\item A política \textit{steal} permite que:

\begin{enumerate}[label=\alph*)]
    \item uma página alterada por transação não confirmada seja escrita no disco.
    \item somente páginas de transações confirmadas sejam escritas.
    \item nenhuma página seja escrita antes do encerramento do SGBD.
    \item apenas o log permaneça na memória.
\end{enumerate}

\item A política \textit{no-steal} determina que:

\begin{enumerate}[label=\alph*)]
    \item páginas alteradas por transações não confirmadas não sejam gravadas no banco.
    \item qualquer página possa ser escrita a qualquer momento.
    \item todo commit exija falha.
    \item o log não seja utilizado.
\end{enumerate}

\item A política \textit{force} determina que:

\begin{enumerate}[label=\alph*)]
    \item todas as atualizações da transação sejam gravadas no banco no momento do commit.
    \item nenhuma atualização seja gravada no commit.
    \item todas as transações sejam abortadas.
    \item o log seja removido.
\end{enumerate}

\item A política \textit{no-force} permite que:

\begin{enumerate}[label=\alph*)]
    \item uma transação confirme antes que todas as páginas alteradas sejam gravadas no banco.
    \item toda página seja obrigatoriamente gravada antes do commit.
    \item nenhuma transação confirme.
    \item o sistema dispense REDO em qualquer situação.
\end{enumerate}

\item A combinação \textit{steal/no-force} normalmente exige:

\begin{enumerate}[label=\alph*)]
    \item UNDO e REDO.
    \item somente UNDO.
    \item somente REDO.
    \item nem UNDO nem REDO.
\end{enumerate}

\item A combinação \textit{no-steal/no-force} normalmente exige:

\begin{enumerate}[label=\alph*)]
    \item somente REDO.
    \item somente UNDO.
    \item UNDO e REDO.
    \item nem UNDO nem REDO.
\end{enumerate}

\item A combinação \textit{steal/force} normalmente exige:

\begin{enumerate}[label=\alph*)]
    \item somente UNDO.
    \item somente REDO.
    \item UNDO e REDO.
    \item nem UNDO nem REDO.
\end{enumerate}

\item A combinação \textit{no-steal/force} normalmente exige:

\begin{enumerate}[label=\alph*)]
    \item nem UNDO nem REDO.
    \item UNDO e REDO.
    \item apenas REDO.
    \item apenas UNDO.
\end{enumerate}

\item Na atualização adiada, as atualizações no banco são realizadas:

\begin{enumerate}[label=\alph*)]
    \item somente depois que a transação alcança seu ponto de efetivação.
    \item imediatamente após cada comando.
    \item somente antes de iniciar a transação.
    \item apenas após uma falha.
\end{enumerate}

\item Antes do commit, uma transação com atualização adiada modifica:

\begin{enumerate}[label=\alph*)]
    \item buffers e registros de log, mas não o banco persistente.
    \item diretamente todas as páginas do banco.
    \item apenas as chaves primárias.
    \item somente o catálogo.
\end{enumerate}

\item Se uma falha ocorre antes do commit em atualização adiada:

\begin{enumerate}[label=\alph*)]
    \item não é necessário UNDO, pois o banco não recebeu as alterações.
    \item é necessário obrigatoriamente UNDO.
    \item a transação deve ser refeita.
    \item o banco inteiro deve ser restaurado de backup.
\end{enumerate}

\item Se uma transação confirmou, mas algumas de suas atualizações ainda não chegaram ao banco:

\begin{enumerate}[label=\alph*)]
    \item pode ser necessário REDO.
    \item pode ser necessário apenas UNDO.
    \item nenhum mecanismo é necessário.
    \item a transação deve ser considerada abortada.
\end{enumerate}

\item O protocolo típico de atualização adiada é classificado como:

\begin{enumerate}[label=\alph*)]
    \item NO-UNDO/REDO.
    \item UNDO/REDO.
    \item UNDO/NO-REDO.
    \item NO-UNDO/NO-REDO.
\end{enumerate}

\item Em NO-UNDO/REDO, uma transação:

\begin{enumerate}[label=\alph*)]
    \item não escreve alterações no banco antes do commit.
    \item sempre escreve no banco antes de registrar no log.
    \item exige UNDO em qualquer falha.
    \item nunca utiliza REDO.
\end{enumerate}

\item Em NO-UNDO/REDO, após uma falha devem ser refeitas:

\begin{enumerate}[label=\alph*)]
    \item transações confirmadas cujos efeitos podem não estar persistidos.
    \item todas as transações incompletas.
    \item apenas transações abortadas.
    \item somente transações que não alteraram dados.
\end{enumerate}

\item Na atualização imediata, uma transação pode:

\begin{enumerate}[label=\alph*)]
    \item atualizar o banco antes de chegar ao commit.
    \item atualizar o banco apenas depois do commit.
    \item executar somente leituras.
    \item dispensar completamente o log.
\end{enumerate}

\item Como a atualização imediata pode alterar o banco antes do commit:

\begin{enumerate}[label=\alph*)]
    \item pode ser necessário desfazer efeitos de transações não confirmadas.
    \item UNDO nunca é necessário.
    \item REDO nunca é necessário.
    \item o log não precisa registrar valores antigos.
\end{enumerate}

\item A forma mais geral da atualização imediata é:

\begin{enumerate}[label=\alph*)]
    \item UNDO/REDO.
    \item NO-UNDO/REDO.
    \item NO-UNDO/NO-REDO.
    \item REDO/NO-COMMIT.
\end{enumerate}

\item Em UNDO/REDO, após uma falha:

\begin{enumerate}[label=\alph*)]
    \item transações não confirmadas são desfeitas e confirmadas podem ser refeitas.
    \item todas as transações são somente refeitas.
    \item todas as transações são somente desfeitas.
    \item nenhuma operação é realizada.
\end{enumerate}

\item O UNDO deve processar os registros normalmente:

\begin{enumerate}[label=\alph*)]
    \item na ordem inversa em que as atualizações ocorreram.
    \item na ordem alfabética dos itens.
    \item sempre do início para o fim.
    \item sem considerar a ordem.
\end{enumerate}

\item O REDO deve processar os registros normalmente:

\begin{enumerate}[label=\alph*)]
    \item na ordem em que foram registrados no log.
    \item sempre na ordem inversa.
    \item de forma aleatória.
    \item somente por nome de tabela.
\end{enumerate}

\item A variação UNDO/NO-REDO ocorre quando:

\begin{enumerate}[label=\alph*)]
    \item alterações não confirmadas podem chegar ao banco, mas todas as alterações confirmadas são forçadas ao disco.
    \item nenhuma alteração não confirmada chega ao banco e alterações confirmadas podem permanecer apenas em memória.
    \item nenhuma alteração chega ao banco.
    \item apenas paginação sombra é utilizada.
\end{enumerate}

\item Em UNDO/NO-REDO, uma transação confirmada não precisa de REDO porque:

\begin{enumerate}[label=\alph*)]
    \item suas alterações são gravadas no banco antes da confirmação ser concluída.
    \item ela nunca atualiza o banco.
    \item seus registros de log são ignorados.
    \item toda transação confirmada é desfeita.
\end{enumerate}

\item Uma técnica classificada como NO-UNDO/NO-REDO é:

\begin{enumerate}[label=\alph*)]
    \item paginação sombra.
    \item atualização imediata \textit{steal/no-force}.
    \item atualização adiada \textit{no-force}.
    \item WAL convencional com UNDO/REDO.
\end{enumerate}

\item Na paginação sombra, a tabela de páginas sombra:

\begin{enumerate}[label=\alph*)]
    \item representa a versão estável anterior do banco.
    \item é sobrescrita a cada pequena modificação.
    \item contém apenas transações abortadas.
    \item não aponta para páginas em disco.
\end{enumerate}

\item Durante a execução com paginação sombra, as alterações são refletidas:

\begin{enumerate}[label=\alph*)]
    \item em uma tabela de páginas corrente e em novas páginas.
    \item diretamente na tabela sombra.
    \item apenas no log.
    \item somente no catálogo.
\end{enumerate}

\item No commit da paginação sombra:

\begin{enumerate}[label=\alph*)]
    \item o ponteiro para a tabela corrente passa a representar a versão oficial.
    \item todas as páginas novas são descartadas.
    \item o banco volta obrigatoriamente à versão antiga.
    \item todas as transações são desfeitas.
\end{enumerate}

\item Em caso de falha antes do commit na paginação sombra:

\begin{enumerate}[label=\alph*)]
    \item a tabela sombra continua apontando para a versão consistente anterior.
    \item é necessário refazer obrigatoriamente todas as alterações.
    \item a versão anterior é perdida.
    \item todo o banco precisa ser apagado.
\end{enumerate}

\item Um checkpoint é usado para:

\begin{enumerate}[label=\alph*)]
    \item limitar a parte do log que precisa ser examinada durante a recuperação.
    \item substituir permanentemente o log.
    \item impedir qualquer atualização.
    \item eliminar backups.
\end{enumerate}

\item Ao criar um checkpoint, o sistema registra:

\begin{enumerate}[label=\alph*)]
    \item um ponto conhecido do estado do sistema e informações sobre transações ativas.
    \item apenas os nomes das tabelas.
    \item somente as chaves estrangeiras.
    \item apenas comandos \texttt{SELECT}.
\end{enumerate}

\item Após uma falha, o checkpoint permite que a recuperação:

\begin{enumerate}[label=\alph*)]
    \item comece a análise a partir de um ponto mais recente do log.
    \item sempre leia o log desde o início da existência do banco.
    \item dispense análise das transações.
    \item confirme todas as transações incompletas.
\end{enumerate}

\item Quanto mais frequentes forem os checkpoints:

\begin{enumerate}[label=\alph*)]
    \item menor pode ser o trabalho de recuperação, mas maior o custo durante execução.
    \item maior será sempre o tempo de recuperação.
    \item menor será o uso de log, sem qualquer custo.
    \item menos consistência o banco terá.
\end{enumerate}

\item Em recuperação multiusuário, o algoritmo precisa considerar:

\begin{enumerate}[label=\alph*)]
    \item transações concorrentes e o protocolo de controle de concorrência.
    \item somente uma transação por vez.
    \item apenas comandos DDL.
    \item ausência de bloqueios.
\end{enumerate}

\item Um protocolo de concorrência estrito facilita a recuperação porque:

\begin{enumerate}[label=\alph*)]
    \item impede leitura e escrita sobre dados produzidos por transações não confirmadas.
    \item permite livre leitura de valores sujos.
    \item elimina o log.
    \item impede qualquer commit.
\end{enumerate}

\item Em recuperação UNDO/REDO com checkpoint, transações ativas no momento da falha devem:

\begin{enumerate}[label=\alph*)]
    \item ser desfeitas.
    \item ser sempre refeitas.
    \item ser consideradas confirmadas.
    \item ser ignoradas.
\end{enumerate}

\item Em recuperação UNDO/REDO com checkpoint, transações confirmadas após o checkpoint podem:

\begin{enumerate}[label=\alph*)]
    \item precisar ser refeitas.
    \item precisar ser sempre desfeitas.
    \item ser consideradas inexistentes.
    \item ser removidas do banco.
\end{enumerate}

\item Uma falha catastrófica pode envolver:

\begin{enumerate}[label=\alph*)]
    \item perda ou dano do disco onde está o banco.
    \item apenas erro de sintaxe em uma consulta.
    \item somente uma leitura suja.
    \item uso de \texttt{ROLLBACK}.
\end{enumerate}

\item Para recuperação de falhas catastróficas é necessário principalmente:

\begin{enumerate}[label=\alph*)]
    \item backup do banco e dos registros necessários.
    \item somente memória cache.
    \item apenas bloqueios compartilhados.
    \item apenas normalização.
\end{enumerate}

\item O backup completo deve incluir normalmente:

\begin{enumerate}[label=\alph*)]
    \item dados do banco e informações de log relevantes.
    \item apenas nomes de usuários.
    \item somente tabelas vazias.
    \item apenas comandos \texttt{SELECT}.
\end{enumerate}

\item O log costuma ser copiado com maior frequência que o banco completo porque:

\begin{enumerate}[label=\alph*)]
    \item permite recuperar alterações realizadas após o último backup.
    \item não contém informações de transações.
    \item sempre ocupa mais espaço que o banco.
    \item substitui definitivamente o backup.
\end{enumerate}

\item Para recuperar após falha catastrófica, uma estratégia é:

\begin{enumerate}[label=\alph*)]
    \item restaurar o backup e reaplicar alterações confirmadas registradas no log.
    \item apagar o backup e manter apenas transações incompletas.
    \item executar UNDO de todas as transações históricas.
    \item criar um banco vazio.
\end{enumerate}

\item A recuperação de banco de dados depende da interação entre:

\begin{enumerate}[label=\alph*)]
    \item gerenciamento de buffers, log, transações e armazenamento persistente.
    \item somente normalização e SQL.
    \item apenas criação de índices.
    \item apenas controle de usuários.
\end{enumerate}

\end{enumerate}

\section{Questões de verdadeiro ou falso}

\begin{enumerate}[resume,label=\textbf{Q\arabic*.}]

\item Julgue as afirmações:

\begin{enumerate}[label=\alph*)]
    \item ( \quad ) A recuperação ajuda a garantir atomicidade.
    \item ( \quad ) A recuperação também ajuda a garantir durabilidade.
    \item ( \quad ) Falhas podem deixar transações incompletas.
    \item ( \quad ) Recuperação não possui relação com consistência.
\end{enumerate}

\item Julgue as afirmações:

\begin{enumerate}[label=\alph*)]
    \item ( \quad ) O log registra informações sobre atualizações.
    \item ( \quad ) A imagem anterior pode ser usada no UNDO.
    \item ( \quad ) A imagem posterior pode ser usada no REDO.
    \item ( \quad ) O log deve existir somente depois da falha.
\end{enumerate}

\item Julgue as afirmações:

\begin{enumerate}[label=\alph*)]
    \item ( \quad ) WAL significa gravação adiantada de log.
    \item ( \quad ) O log deve chegar ao disco antes da página alterada.
    \item ( \quad ) Registros de log de uma transação devem estar persistidos antes do commit.
    \item ( \quad ) WAL determina que a página seja escrita antes do log.
\end{enumerate}

\item Julgue as afirmações:

\begin{enumerate}[label=\alph*)]
    \item ( \quad ) UNDO restaura valores anteriores.
    \item ( \quad ) REDO reaplica valores posteriores.
    \item ( \quad ) UNDO costuma percorrer operações na ordem inversa.
    \item ( \quad ) REDO é utilizado apenas para transações abortadas.
\end{enumerate}

\item Julgue as afirmações:

\begin{enumerate}[label=\alph*)]
    \item ( \quad ) \textit{Steal} permite gravar página de transação não confirmada.
    \item ( \quad ) \textit{No-steal} evita que efeitos não confirmados cheguem ao banco.
    \item ( \quad ) \textit{Force} grava atualizações no commit.
    \item ( \quad ) \textit{No-force} elimina sempre a necessidade de REDO.
\end{enumerate}

\item Julgue as afirmações:

\begin{enumerate}[label=\alph*)]
    \item ( \quad ) Atualização adiada não altera o banco antes do commit.
    \item ( \quad ) Atualização adiada típica é NO-UNDO/REDO.
    \item ( \quad ) Transações incompletas em atualização adiada precisam sempre de UNDO.
    \item ( \quad ) Transações confirmadas podem precisar de REDO.
\end{enumerate}

\item Julgue as afirmações:

\begin{enumerate}[label=\alph*)]
    \item ( \quad ) Atualização imediata pode modificar o banco antes do commit.
    \item ( \quad ) UNDO/REDO é uma técnica de atualização imediata.
    \item ( \quad ) Em UNDO/REDO, transações ativas são desfeitas.
    \item ( \quad ) Em UNDO/REDO, transações confirmadas nunca precisam de REDO.
\end{enumerate}

\item Julgue as afirmações:

\begin{enumerate}[label=\alph*)]
    \item ( \quad ) UNDO/NO-REDO pode usar política \textit{steal/force}.
    \item ( \quad ) NO-UNDO/REDO pode usar política \textit{no-steal/no-force}.
    \item ( \quad ) NO-UNDO/NO-REDO pode ser obtido com \textit{no-steal/force}.
    \item ( \quad ) \textit{Steal/no-force} dispensa UNDO e REDO.
\end{enumerate}

\item Julgue as afirmações:

\begin{enumerate}[label=\alph*)]
    \item ( \quad ) Paginação sombra mantém uma versão estável anterior.
    \item ( \quad ) Alterações são feitas em páginas novas.
    \item ( \quad ) O commit pode trocar o ponteiro da tabela de páginas.
    \item ( \quad ) Paginação sombra básica exige sempre UNDO e REDO.
\end{enumerate}

\item Julgue as afirmações:

\begin{enumerate}[label=\alph*)]
    \item ( \quad ) Checkpoints reduzem o trecho do log analisado na recuperação.
    \item ( \quad ) Checkpoints podem registrar transações ativas.
    \item ( \quad ) Checkpoints frequentes podem aumentar o custo normal de execução.
    \item ( \quad ) Checkpoints tornam o log desnecessário.
\end{enumerate}

\item Julgue as afirmações:

\begin{enumerate}[label=\alph*)]
    \item ( \quad ) Recuperação multiusuário depende do controle de concorrência.
    \item ( \quad ) Escalonamentos estritos facilitam recuperação.
    \item ( \quad ) Em escalonamento estrito, valores não confirmados ficam livremente disponíveis.
    \item ( \quad ) Deadlocks podem causar abortos que exigem UNDO.
\end{enumerate}

\item Julgue as afirmações:

\begin{enumerate}[label=\alph*)]
    \item ( \quad ) Falha catastrófica pode destruir o armazenamento principal do banco.
    \item ( \quad ) Backup é necessário para falhas catastróficas.
    \item ( \quad ) O log pode permitir reaplicar alterações após o último backup.
    \item ( \quad ) Backup completo nunca inclui dados do banco.
\end{enumerate}

\end{enumerate}

\section{Questões de aplicação objetiva}

\begin{enumerate}[resume,label=\textbf{Q\arabic*.}]

\item Uma transação $T_1$ alterou o saldo de R\$500 para R\$300, mas sofreu falha antes do commit. O banco já contém R\$300 e o log possui os dois valores. Qual ação é necessária?

\begin{enumerate}[label=\alph*)]
    \item REDO para gravar R\$300 novamente.
    \item UNDO para restaurar R\$500.
    \item Nenhuma ação.
    \item Confirmar automaticamente a transação.
\end{enumerate}

\item Uma transação confirmou uma alteração de R\$500 para R\$300, mas a página modificada ainda não havia sido gravada no disco quando ocorreu uma falha. Qual ação pode ser necessária?

\begin{enumerate}[label=\alph*)]
    \item UNDO para restaurar R\$500.
    \item REDO para reaplicar R\$300.
    \item Abortar definitivamente todas as transações.
    \item Excluir o log.
\end{enumerate}

\item Em um sistema de atualização adiada, $T_1$ falha antes do commit. O que deve acontecer?

\begin{enumerate}[label=\alph*)]
    \item Deve ser executado UNDO no banco.
    \item Nenhum UNDO é necessário, pois suas alterações não chegaram ao banco.
    \item Todas as alterações devem ser refeitas.
    \item O backup completo deve ser restaurado.
\end{enumerate}

\item Em um sistema NO-UNDO/REDO, $T_1$ confirmou e $T_2$ estava ativa no momento da falha. A recuperação deve:

\begin{enumerate}[label=\alph*)]
    \item refazer $T_1$ se necessário e ignorar os efeitos de $T_2$ no banco.
    \item desfazer $T_1$ e refazer $T_2$.
    \item refazer ambas obrigatoriamente.
    \item desfazer ambas.
\end{enumerate}

\item Em um sistema UNDO/REDO, $T_1$ confirmou e $T_2$ estava ativa na falha. A recuperação deve:

\begin{enumerate}[label=\alph*)]
    \item executar UNDO de $T_1$ e REDO de $T_2$.
    \item executar REDO de $T_1$ e UNDO de $T_2$.
    \item executar UNDO de ambas.
    \item ignorar ambas.
\end{enumerate}

\item Um sistema permite que páginas de transações não confirmadas sejam gravadas no banco e também permite commit antes de todas as páginas chegarem ao disco. Qual algoritmo é necessário?

\begin{enumerate}[label=\alph*)]
    \item NO-UNDO/NO-REDO.
    \item NO-UNDO/REDO.
    \item UNDO/NO-REDO.
    \item UNDO/REDO.
\end{enumerate}

\item Um sistema não permite gravar páginas de transações não confirmadas, mas permite commit antes da gravação de todas as páginas. Qual algoritmo é adequado?

\begin{enumerate}[label=\alph*)]
    \item NO-UNDO/REDO.
    \item UNDO/REDO.
    \item UNDO/NO-REDO.
    \item NO-UNDO/NO-REDO.
\end{enumerate}

\item Um sistema pode gravar páginas de transações não confirmadas, mas força todas as páginas da transação no commit. Qual algoritmo é adequado?

\begin{enumerate}[label=\alph*)]
    \item NO-UNDO/REDO.
    \item UNDO/NO-REDO.
    \item UNDO/REDO.
    \item NO-UNDO/NO-REDO.
\end{enumerate}

\item Um sistema não grava páginas não confirmadas e força todas as páginas no commit. Qual algoritmo pode ser utilizado?

\begin{enumerate}[label=\alph*)]
    \item UNDO/REDO.
    \item NO-UNDO/REDO.
    \item UNDO/NO-REDO.
    \item NO-UNDO/NO-REDO.
\end{enumerate}

\item Considere o log:

\[
\langle start,T_1\rangle
\]
\[
\langle T_1,X,10,20\rangle
\]
\[
\langle commit,T_1\rangle
\]
\[
\langle start,T_2\rangle
\]
\[
\langle T_2,Y,30,50\rangle
\]

O sistema falha logo depois. Em UNDO/REDO, deve-se:

\begin{enumerate}[label=\alph*)]
    \item desfazer $T_1$ e refazer $T_2$.
    \item refazer $T_1$ e desfazer $T_2$.
    \item refazer ambas.
    \item desfazer ambas.
\end{enumerate}

\item Para desfazer a atualização de $Y$ de 30 para 50, o sistema deve gravar em $Y$:

\begin{enumerate}[label=\alph*)]
    \item 50.
    \item 30.
    \item 20.
    \item 10.
\end{enumerate}

\item Para refazer a atualização de $X$ de 10 para 20, o sistema deve gravar em $X$:

\begin{enumerate}[label=\alph*)]
    \item 10.
    \item 20.
    \item 30.
    \item 50.
\end{enumerate}

\item Uma página alterada é gravada no disco antes de seu registro de log. Logo depois ocorre uma falha. Qual princípio foi violado?

\begin{enumerate}[label=\alph*)]
    \item Two-Phase Locking.
    \item Write-Ahead Logging.
    \item Primeira Forma Normal.
    \item Atomicidade de esquema.
\end{enumerate}

\item O último backup completo foi realizado no domingo. Na quarta-feira, o disco foi perdido, mas há cópias do log até o momento da falha. A recuperação adequada é:

\begin{enumerate}[label=\alph*)]
    \item restaurar o backup de domingo e reaplicar transações confirmadas do log.
    \item utilizar apenas \texttt{ROLLBACK}.
    \item descartar o backup e criar um banco vazio.
    \item desfazer todas as transações desde a criação do banco.
\end{enumerate}

\item Um banco utiliza paginação sombra e falha antes do commit. A recuperação deve usar:

\begin{enumerate}[label=\alph*)]
    \item a tabela de páginas sombra que aponta para a versão consistente anterior.
    \item UNDO obrigatório de todas as páginas.
    \item REDO obrigatório de todas as páginas.
    \item apenas o último comando SQL executado.
\end{enumerate}

\end{enumerate}

\newpage

\section*{Gabarito}

\begin{multicols}{2}

\begin{enumerate}[label=\textbf{Q\arabic*.}]
    \item b
    \item a
    \item a
    \item a
    \item a
    \item a
    \item a
    \item a
    \item a
    \item b
    \item a
    \item b
    \item a
    \item a
    \item a
    \item a
    \item a
    \item a
    \item a
    \item a
    \item a
    \item a
    \item a
    \item a
    \item a
    \item a
    \item a
    \item a
    \item a
    \item a
    \item a
    \item a
    \item a
    \item a
    \item a
    \item a
    \item a
    \item a
    \item a
    \item a
    \item a
    \item a
    \item a
    \item a
    \item a
    \item a
    \item a
    \item a
    \item a
    \item a
    \item a
    \item a
    \item a
    \item a
    \item a
    \item a
    \item a
    \item a
    \item a
    \item a
    \item a
    \item a
    \item a
    \item V, V, V, F
    \item V, V, V, F
    \item V, V, V, F
    \item V, V, V, F
    \item V, V, V, F
    \item V, V, F, V
    \item V, V, V, F
    \item V, V, V, F
    \item V, V, V, F
    \item V, V, V, F
    \item V, V, F, V
    \item V, V, V, F
    \item b
    \item b
    \item b
    \item a
    \item b
    \item d
    \item a
    \item b
    \item d
    \item b
    \item b
    \item b
    \item b
    \item a
    \item a
\end{enumerate}

\end{multicols}

\end{document}
